% =============================================================
% File: chapters/ta/chapter-3/analisis-masalah.tex
% =============================================================

\section{Analisis Masalah}

\subsection{Analisis Sistem DecMed}
\label{analisis-decmed}

Sebagaimana dijelaskan pada Subbab~\ref{sec:penelitian-terkait}, DecMed adalah kerangka kerja manajemen RME terdesentralisasi dengan menggunakan DLT IOTA sebagai penyimpanan \textit{on-chain}, IPFS sebagai penyimpanan \textit{off-chain}, serta PRE untuk mendukung distribusi akses data yang aman \cite{DecMed}. Arsitektur ini sudah memberikan fondasi untuk menjaga integritas data dan mengurangi ketergantungan pada penyimpanan terpusat, sebagaimana telah diuraikan pada Subbab~\ref{subsec:dlt-umum}, \ref{subsec:iota}, dan \ref{subsec:pola-layanan-decmed}.

Meskipun arsitektur tersebut telah berhasil meningkatkan aspek keamanan penyimpanan data, DecMed lebih berfokus pada pencegahan serangan dan perlindungan data melalui kontrol akses. Aspek audit log sebagai mekanisme untuk menghasilkan, menyimpan, dan mengelola bukti digital selama proses investigasi insiden keamanan masih belum menjadi fokus pembahasan, sebagaimana ditunjukkan pada Subbab~\ref{sec:penelitian-terkait}, di mana DecMed baru mengklaim ketersediaan \textit{audit trail} secara umum tanpa merinci komponen-komponen audit log yang telah diuraikan pada Subbab~\ref{subsec:komponen-audit-log}.

\subsubsection{Implementasi Audit Log pada DecMed}
\label{subsubsec:implementasi-audit-log}

Dalam operasionalnya, implementasi DecMed sebenarnya telah menyediakan beberapa bentuk pencatatan aktivitas sistem. Pencatatan ini dilakukan melalui dua komponen utama, yaitu \textit{PatientAccessLog} yang dikelola di dalam sistem dan \textit{transaction block} bawaan pada jaringan IOTA. Kedua komponen tersebut telah mendokumentasikan sebagian aktivitas yang terjadi selama sistem beroperasi. Meskipun demikian, apabila dinilai terhadap karakteristik audit log yang baik pada Subbab~\ref{subsec:karakteristik-audit-log}, kedua upaya ini belum membentuk sebuah audit log yang utuh dan mampu memenuhi kebutuhan audit maupun investigasi insiden keamanan karena masing-masing memiliki keterbatasan yang mendasar.

\subsubsection{Keterbatasan \textit{PatientAccessLog}}
\label{subsubsec:keterbatasan-patientaccesslog}

\textit{PatientAccessLog} digunakan untuk mencatat aktivitas pemberian hak akses dari pasien kepada tenaga medis. Informasi yang tersimpan telah mencakup atribut penting seperti identitas pasien, identitas tenaga medis, waktu pemberian akses (\textit{timestamp}), serta informasi lain yang berkaitan dengan hak akses tersebut.

Namun demikian, \textit{PatientAccessLog} memiliki dua keterbatasan kritis sebagai sebuah komponen audit log. Pertama, cakupan pencatatannya masih sangat terbatas pada proses pemberian hak akses. Berbagai aktivitas lain yang memiliki nilai audit tinggi, seperti autentikasi, otorisasi, perubahan konfigurasi, maupun interaksi kueri terhadap data rekam medis belum terdokumentasi. Hal ini bertentangan dengan cakupan minimal peristiwa pemicu (\textit{trigger events}) yang telah diuraikan pada Subbab~\ref{subsubsec:penentuan-event}. Kedua, mekanisme ini rentan terhadap perubahan sehingga tidak memenuhi karakteristik kekekalan data (\textit{append-only}). \textit{PatientAccessLog} dikaitkan secara langsung dengan fungsionalitas pencabutan hak akses (\textit{revoke access}). Ketika hak akses dicabut, informasi pada entri log yang sama akan diperbarui dengan menimpa (\textit{overwrite}) status akses sebelumnya menjadi \textit{revoked}. Pendekatan ini melanggar karakteristik integritas audit log pada Subbab~\ref{subsec:karakteristik-audit-log}, yang mensyaratkan bahwa catatan historis tidak boleh berubah dan aktivitas baru seharusnya dicatat sebagai \textit{audit record} terpisah sesuai definisi pada Subbab~\ref{subsec:definisi-audit-log}.

\subsubsection{Keterbatasan \textit{Transaction Block} IOTA}
\label{subsubsec:keterbatasan-tx-block}

Selain pencatatan internal, DecMed memanfaatkan \textit{transaction block} pada jaringan IOTA sebagai bukti transaksi. Setiap pemanggilan fungsi \textit{smart contract} akan menghasilkan \textit{transaction block} yang tersimpan secara terdistribusi. Mekanisme ini memberikan jaminan integritas yang tinggi karena data bersifat \textit{immutable} dan sulit dimodifikasi, sejalan dengan karakteristik DLT pada Subbab~\ref{subsec:dlt-umum}.

Meskipun unggul dalam hal integritas, pemanfaatan \textit{transaction block} IOTA sebagai sumber audit log masih memiliki beberapa kelemahan utama. Kelemahan pertama adalah inkonsistensi format data, di mana \textit{audit record} yang dihasilkan belum terstandarisasi karena struktur data transaksi sangat bervariasi bergantung pada jenis fungsi atau operasi spesifik yang sedang dieksekusi oleh sistem. Selanjutnya, terdapat keterbatasan pada variasi \textit{event} yang dicatat. Cakupan peristiwa yang terekam sangat sempit karena komponen ini hanya mendokumentasikan eksekusi operasi \textit{move call} pada \textit{smart contract}, sehingga berbagai interaksi sistem yang tidak memicu pemanggilan fungsi tersebut akan luput dari pencatatan. Kelemahan terakhir berkaitan dengan permasalahan retensi data. Ketersediaan bukti digital dalam jangka panjang tidak dapat dijamin secara pasti karena data transaksi pada \textit{node} IOTA rentan terhapus akibat mekanisme \textit{data pruning} yang rutin dilakukan untuk membebaskan ruang penyimpanan \textit{node}. Kondisi ini membuat pemanfaatan \textit{transaction block} belum memenuhi persyaratan retensi komponen penyimpanan \textit{event} sebagaimana diuraikan pada Subbab~\ref{subsubsec:penyimpanan-event}.

\subsubsection{Keterbatasan Pencatatan Audit Log yang Tersebar pada DecMed}
\label{subsubsec:keterbatasan-tersebar}

Implementasi pencatatan aktivitas pada DecMed saat ini masih bersifat tersebar dan terisolasi pada masing-masing komponen internal tanpa adanya pengelola terpusat. Kondisi arsitektur pencatatan yang terfragmentasi ini menimbulkan sejumlah keterbatasan kritis yang menyulitkan pelaksanaan pembuktian dan investigasi digital ketika terjadi insiden keamanan.

Keterbatasan pertama terletak pada fragmentasi format dan ketiadaan standar log. Karena setiap komponen penghasil kejadian (\textit{event source}) memproduksi log secara independen, skema dan atribut data yang dihasilkan menjadi sangat heterogen. Ketiadaan standarisasi ini secara langsung menyulitkan proses korelasi dan analisis silang (\textit{cross-analysis}) antar-komponen secara otomatis. Selain itu, sistem pencatatan yang tersebar mengandalkan pewaktu lokal dari masing-masing lingkungan (\textit{environment}) operasional. Perbedaan sumber pewaktu ini memunculkan risiko kejanggalan waktu (\textit{time-skew}) antar-log, sehingga rekonstruksi kronologi kejadian yang presisi menjadi sangat sulit untuk dilakukan.

Keterbatasan selanjutnya adalah kesulitan dalam melakukan agregasi bukti digital secara menyeluruh. Tidak tersedianya layanan pengumpul terpusat (\textit{centralized collector}) menyebabkan pemantauan aktivitas sistem tidak dapat dilakukan secara holistik, dan pengumpulan bukti harus ditarik secara parsial dari berbagai tempat yang berbeda. Hal ini diperparah oleh ketiadaan strategi retensi data yang seragam. Komponen yang tersebar memiliki mekanisme dan batas kapasitas penyimpanan masing-masing. Sebagian catatan audit mungkin tersimpan di basis data lokal tanpa jaminan perlindungan terhadap manipulasi, sementara catatan lainnya yang berada pada jaringan DLT rentan mengalami penghapusan permanen akibat mekanisme pembersihan (\textit{pruning}) secara berkala.